\chapter{Systen Design}
\label{chap:design}
In this chapter, we lay out the foundations of the problem at hand, and how did we design a solution to approach it. We begin with our assumptions about the environment, and then detail out individual design goals that were chosen because they correlate with an assumption.
\section{Environment Assumptions}
\label{sec:ass}
The operating environment for GSN is versatile: from disaster-struck warzones that lack physical access to the Internet because its operating infrastructure was damaged, to politically oppresive regimes that use various Internet monitoring and controling techniques to oppress and cull resistance to them, providing secure communications in times of political unrest, and providing an end-to-end secure messaging platform for privacy-aware individuals.
Due to the various possible settings in which it can be used, the system was designed with letting the user choice of the right security settings for them, as provides a tradeoff between usability and security and privacy, while the default settings for all features are set to the most-cautious, unless chosen otherwise.
The first environment that we decided to operate in is disaster zones, which usually have damaged wide-area Internet infrastructure and backhaul connectivity due to the infrastructure's delicate nature~\cite{Legendre2011,Liu2026,Hasan2025}. In those scenarios, only a small majority of nodes (e.g. those with satellite links like Starlink plates) retain access to the global Internet~\cite{Liu2026}. Due to the terrain fragmentation and partial Internet backhaul connectivity, research propose various OppNet and DTN routing strategies, most notably epidemic routing, to pass critical messages directly to one another until reaching a node that still has a connection. To maximize device life and minimize battery drain, we thus decided to focus on BLE as a main transport medium, and to develop a mesh strategy to maximize the transmission range and overall reachability, using other peers to forward messages.
Other then disaster zones, another motivating environment is the issue of providing secure communication in politically oppressive regimes, that still have a functioning Internet infrastructure. The underlying transport layer protocol choice was shaped by this design space: TCP is heavily censored in Iran, by using stateful inspection to detect and block banned domains~\cite{Aceto2026}, while UDP is heavily restrained due to being associated with VPN traffic and the QUIC protocol, which is used by major Google platforms like Search, YouTube, and Gmail, and Meta CSNs like Facebook and Instagram. While network anomalities targeting the QUIC protocol were detecting during the January 2026 Internet outage in Iran~\cite{Aceto2026}, researchers noted that 4.4\% of global Internet users cannot use QUIC because UDP is completely blocked, however most of those users were found to be corporate network users that are behind restrictive enterprise firewalls, and no single ISP was found to completely block QUIC or UDP~\cite{Langley2017}.
As mentioned in the previous chapter, we chose UDP as our underlying transport protocol for the IP Internet relation, using it to establish secure communication channels with remote peers. To establish those channels, a hole needs to be punched in each user's firewall to let data packets in from their counterpart. For that, a communication channel needs to be establish to be used as a signaling channel, to coordinate hole punching. Initially using the BLE transport to directly coordinate UDP hole-punching, GNL quickly expanded into facilitating hole-punching through existing network connections with mutual online friends (i.e. edge rendezvous points), which uses the notion of propagating trust along the user's social graph to facilitate connections for friends. Additionally, GNL also supports cloud rendezvous points, which are servers that run the rendezvous points, that any GSN user can run on a server with a static IP address and a permissive firewall policy, that allows external incoming connections. Using cloud rendezvous point allows users to reach remote friends without having any overlap in their GSN social graph or having encountered them physically, as rendezvous points allow the creation and redemption of invite links, which are essentially ephemeral, time-constrained, usage-limited expiring tokens that allow the redeemer to cold-call the link creator, letting the redeemer submit a friend request through one of the creator's rendezvous points (both edge and cloud).
Pivoting from a traditional client-server architecture to a p2p architecture was clear from the beginning, as the p2p architecture is much more resistant to censorship and control, while the client-server architecture is also more subsceptible to distruptions.
Finally, to mitigate online monitoring, we decided to enforce end-to-end encryption using a secure Noise handshake for an encrypted session establishment, rendering all transported text undecipherable ciphertext to an adversary that has intercepted a message.

\section{Design Goals}
\label{sec:des}
This section presents design goals and the philosophical foundation for implementation of the network layer. These design goals were derived from the assumptions in section\ref{sec:ass}, and are reflected throughout the system.

\textbf{Local over remote}: The system should prefer local over remote communication.
When possible, GNL must use BLE over IP communication to deliver messages: this means, when a user pair has a link with both the BLE and IP protocols, the layer will choose BLE as a carrier medium for a message. This is because BLE involves less external components and exposes less information for third parties.

\textbf{Internet over long distance}: When not in close range, the system must use available Internet intrastructure to enable IP communication.
When a user pair is befriended, GNL should use IP to convey a message between two remote peers. This means the system will choose to relay a message over the IP medium rather than wait for an opportunistic delivery with the BLE mesh.

\textbf{Reduce attack interface}: The system should expose the minimum amount of required information to strangers in order to facilitate communication.
In order to facilitate mesh communication, the cleartext part of the protocol header should be reduced as much as possible. This will harden the system's security guarantees by making it hard to perform profiling attacks, and decrypt ciphertext.

\textbf{Store-carry-forward}: The systen should cache packets for unavailable recipients until a possible opportunistic delivery attempt can be made.
Digressing from traditional IP routing, the system will try opportunistic delivery of messages for offline peers, queueing messages until a suitable carrier can be found.

\textbf{User choice}: The system must let the user decide on various secure privacy-critical settings.
Due to being used in various security profiles affected by different usage scenarios, the system lets the user decide on the usability tradeoff themselves, whether to use sensitive options like a cloud RV point (which is sensitive to control and outages), invite links, IP communication, and more. These options should be configurable in a settings menu.

\section{Architecture}
\label{sec:arch}

A demo chat app deployed on phones, which serves as the foundation to the blocklace-based future grassroots social network: argue why exactly these functions are needed.
\section{System Architecture}

\resizebox{!}{0.33\textheight}{ % Resizes height to 1/3 of the text height, aspect ratio preserved
        \begin{tikzpicture}[
            stack/.style={
                draw, 
                thick, 
                minimum height=1.2cm, 
                font=\sffamily\Huge,
                inner sep=10pt
            },
            arrow/.style={
                ->, 
                >=stealth, 
                thick
            }
        ]

            % --- Layer 1: Bottom (BLE, Cellular, WLAN) ---
            \node[stack, minimum width=3cm] (ble) {BLE};
            \node[stack, minimum width=3cm, anchor=west] (cellular) at ($(ble.east)+(0.5cm,0)$) {Cellular};
            \node[stack, minimum width=3cm, anchor=west] (wlan) at ($(cellular.east)+(0.5cm,0)$) {WLAN};

            % --- Layer 2: IP (Full Width) ---
            \node[stack, minimum width=10cm, anchor=south west] (ip) at ($(ble.north west)+(0,0.8cm)$) {IP};

            % --- Layer 3: UDP with TLS to the right ---
            \node[stack, minimum width=7cm, anchor=south west] (udp) at ($(ip.north west)+(0,0.8cm)$) {UDP};
            \node[stack, minimum width=2.5cm, anchor=south east] (noise) at ($(ip.north east)+(0,0.8cm)$) {Noise};

            % --- Layer 4: Top (GSN - Full Width) ---
            \node[stack, minimum width=10cm, anchor=south west] (gsn) at ($(udp.north west)+(0,0.8cm)$) {GSN};

            % --- Arrows (Downward) ---
            % From GSN to UDP
            \draw[arrow] (gsn.south -| udp.north) -- (udp.north);
            \draw[arrow] (gsn.south -| noise.north) -- (noise.north);

            % From UDP to IP
            \draw[arrow] (udp.south) -- (udp.south |- ip.north);

            % From IP to BLE, Cellular, WLAN
            \draw[arrow] (ip.south -| ble.north) -- (ble.north);
            \draw[arrow] (ip.south -| cellular.north) -- (cellular.north);
            \draw[arrow] (ip.south -| wlan.north) -- (wlan.north);

        \end{tikzpicture}
    }

% ```
% ┌─────────────────────────────────────────────────────────────────────────┐
% │  Application layer  (GSG / madGLP runtime; here: demo chat app)         │
% │  • Consumes the GrassrootsNetwork façade + reads the Redux store        │
% │  • Supplies opaque payloads (Blocks); renders peers, chats, settings    │
% ├─────────────────────────────────────────────────────────────────────────┤
% │  Coordinator  —  GrassrootsNetwork  (lib/src/grassroots_network.dart)   │
% │  • Wires transports, protocol, sessions, signaling, and the store       │
% │  • Owns the send/receive pipeline, the outbound queue, retry logic      │
% ├─────────────────────────────────────────────────────────────────────────┤
% │  Protocol & security                                                    │
% │  • Packet codec        (models/packet.dart, protocol/protocol_handler)  │
% │  • Fragmentation       (protocol/fragment_handler.dart)                 │
% │  • Noise XX sessions   (session/noise_session_manager.dart)             │
% │  • Dedup / routing     (routing/message_router.dart, mesh/bloom_filter) │
% │  • Signaling / NAT     (signaling/*, transport/hole_punch_service.dart) │
% ├─────────────────────────────────────────────────────────────────────────┤
% │  Transport layer  —  TransportService interface                         │
% │  • BleTransportService  (dual-role GATT, grassroots_bluetooth_layer)    │
% │  • UdpTransportService  (UDX over dual-stack UDP)                       │
% ├─────────────────────────────────────────────────────────────────────────┤
% │  Platform / hardware  (iOS & Android BLE radios, OS UDP sockets)        │
% └─────────────────────────────────────────────────────────────────────────┘

%            ┌──────────────────────────────────────────────┐
%            │        Redux store (AppState)                │  ← cross-cutting:
%            │  peers · transports · messages · friendships │    a strict projection
%            │  · settings · signaling                      │    of transport facts
%            └──────────────────────────────────────────────┘
% ```

\begin{itemize}
    \item Ed25519 identity
    \item Two transport protocols, one interface
    \item Message types: announce, message, friendship
    \item Friend vs. Stranger
    \item Trust Levels: cold calls
    \item Facilitation opt-in: edge RV points
    \item Always-on cloud RV points: optionally use RVs to facilitate connections
    \item Signaling over BT or mutual friends: hole-punching 
    \item Simple messaging: text and media
    \item Blocks: ACK, NACK
\end{itemize}